\section{Sextic promotion and two period-two scales}
\label{sec:sextic}

\subsection{Sixth-order interpolating Hamiltonian}

The next term is
\begin{equation}
\boxed{
H_6
=
dQ^6+eQ^4P^2+fQ^2P^4+gP^6.
}
\label{eq:H6}
\end{equation}
Define
\begin{equation}
\boxed{
G=d+e+f+g.
}
\label{eq:G-def}
\end{equation}

The detuned truncated normal form is
\begin{equation}
\boxed{
K_\delta
=
-\delta I+H_4+H_6
+
\bigO(\delta^2I,\delta I^2,I^4).
}
\label{eq:K-detuned}
\end{equation}

\subsection{Axial hyperbolic branches}

On the \(Q\)-axis,
\[
K_Q
=
-\delta I+4AI^2+8dI^3+\cdots.
\]
Hence
\[
I_Q
=
\frac{\delta}{8A}
-
\frac{3d}{64A^3}\delta^2
+
\bigO(\delta^3),
\]
and
\begin{equation}
\boxed{
\mcJ_Q
=
-\frac{\delta^2}{16A}
+
\frac{d}{64A^3}\delta^3
+
\bigO(\delta^4).
}
\label{eq:JQ}
\end{equation}

Similarly,
\begin{equation}
\boxed{
\mcJ_P
=
-\frac{\delta^2}{16A}
+
\frac{g}{64A^3}\delta^3
+
\bigO(\delta^4).
}
\label{eq:JP}
\end{equation}
Therefore
\begin{equation}
\boxed{
|\mcJ_H|\asymp\delta^2.
}
\label{eq:axial-action}
\end{equation}

The physical second-return exponents satisfy
\begin{equation}
\boxed{
\kappa_{F,Q}
=
2\rho\delta
+
\frac{3d-e}{4A^2\rho}\delta^2
+
\bigO(\delta^3),
}
\label{eq:kappaQ}
\end{equation}
and
\begin{equation}
\boxed{
\kappa_{F,P}
=
2\rho\delta
+
\frac{3g-f}{4A^2\rho}\delta^2
+
\bigO(\delta^3).
}
\label{eq:kappaP}
\end{equation}

\subsection{Near-diagonal scalar reduction}

Along the exact diagonal, \(Q^2=P^2=I\), so the leading scalar model is
\begin{equation}
\boxed{
K_E(I)
=
-\delta I+aI^2+GI^3,
\qquad
a=A\sigma<0.
}
\label{eq:diagonal-scalar}
\end{equation}

The positive stationary solution is
\begin{equation}
\boxed{
I_E
=
\frac{|a|+\sqrt{a^2+3G\delta}}{3G}.
}
\label{eq:IE}
\end{equation}

Define
\begin{equation}
\boxed{
\deltax=\frac{a^2}{G}.
}
\label{eq:crossover-scale}
\end{equation}

For \(\delta\gg\deltax\),
\[
I_E\sim\sqrt{\frac{\delta}{3G}},
\qquad
r_E\asymp\delta^{1/4}.
\]

The reduced action is
\[
\mcJ_E=-aI_E^2-2GI_E^3,
\]
so
\begin{equation}
\boxed{
\mcJ_E
\sim
-\frac{2}{3^{3/2}\sqrt G}\delta^{3/2}.
}
\label{eq:elliptic-action-scaling}
\end{equation}

\subsection{Two-scale periodic-orbit geometry}

The axial scale is
\[
r_H\asymp\delta^{1/2},
\]
whereas the near-diagonal scale is
\[
r_E\asymp\delta^{1/4}.
\]
Hence
\[
\frac{r_E}{r_H}\asymp\delta^{-1/4}.
\]

Invariantly,
\[
\boxed{
|\mcJ_H|\asymp\delta^2,
\qquad
|\mcJ_E|\asymp\delta^{3/2}.
}
\]

\subsection{Breakdown of the outer scaling}

Because \(a<0\) is fixed, the sextic power law cannot persist to \(\delta=0\). Instead,
\[
I_E\to I_0^{\rm scalar}
=
\frac{2|a|}{3G},
\]
and
\[
\mcJ_E\to
\mcJ_0^{\rm scalar}
=
-\frac{4|a|^3}{27G^2}.
\]
The next section describes this crossover.
